\documentclass[10pt,a4paper]{amsart}
\usepackage{amsmath,amssymb,amsthm,graphicx}
\usepackage{verbatim}
\newenvironment{reference}{\comment}{\endcomment}
%\newenvironment{reference}{}{}
\newenvironment{slogan}{\comment}{\endcomment}
\newenvironment{history}{\comment}{\endcomment}
% For commutative diagrams we use Xy-pic
\usepackage[all]{xy}

% We use 2cell for 2-commutative diagrams.
\xyoption{2cell}
% Theorem environments.
%
\theoremstyle{plain}
\newtheorem{theorem}[subsection]{Theorem}
\newtheorem{proposition}[subsection]{Proposition}
\newtheorem{lemma}[subsection]{Lemma}

\theoremstyle{definition}
\newtheorem{definition}[subsection]{Definition}
\newtheorem{example}[subsection]{Example}
\newtheorem{exercise}[subsection]{Exercise}
\newtheorem{situation}[subsection]{Situation}

\theoremstyle{remark}
\newtheorem{remark}[subsection]{Remark}
\newtheorem{remarks}[subsection]{Remarks}

\numberwithin{equation}{subsection}

% Macros
%
\def\lim{\mathop{\mathrm{lim}}\nolimits}
\def\colim{\mathop{\mathrm{colim}}\nolimits}
\def\Spec{\mathop{\mathrm{Spec}}}
\def\Hom{\mathop{\mathrm{Hom}}\nolimits}
\def\Ext{\mathop{\mathrm{Ext}}\nolimits}
\def\SheafHom{\mathop{\mathcal{H}\!\mathit{om}}\nolimits}
\def\SheafExt{\mathop{\mathcal{E}\!\mathit{xt}}\nolimits}
\def\Sch{\mathit{Sch}}
\def\Mor{\mathop{\mathrm{Mor}}\nolimits}
\def\Ob{\mathop{\mathrm{Ob}}\nolimits}
\def\Sh{\mathop{\mathit{Sh}}\nolimits}
\def\NL{\mathop{N\!L}\nolimits}
\def\CH{\mathop{\mathrm{CH}}\nolimits}
\def\proetale{{pro\text{-}\acute{e}tale}}
\def\etale{{\acute{e}tale}}
\def\QCoh{\mathit{QCoh}}
\def\Ker{\mathop{\mathrm{Ker}}}
\def\Im{\mathop{\mathrm{Im}}}
\def\Coker{\mathop{\mathrm{Coker}}}
\def\Coim{\mathop{\mathrm{Coim}}}

% Boxtimes
%
\DeclareMathSymbol{\boxtimes}{\mathbin}{AMSa}{"02}

%
% Macros for moduli stacks/spaces
%
\def\QCohstack{\mathcal{QC}\!\mathit{oh}}
\def\Cohstack{\mathcal{C}\!\mathit{oh}}
\def\Spacesstack{\mathcal{S}\!\mathit{paces}}
\def\Quotfunctor{\mathrm{Quot}}
\def\Hilbfunctor{\mathrm{Hilb}}
\def\Curvesstack{\mathcal{C}\!\mathit{urves}}
\def\Polarizedstack{\mathcal{P}\!\mathit{olarized}}
\def\Complexesstack{\mathcal{C}\!\mathit{omplexes}}
% \Pic is the operator that assigns to X its picard group, usage \Pic(X)
% \Picardstack_{X/B} denotes the Picard stack of X over B
% \Picardfunctor_{X/B} denotes the Picard functor of X over B
\def\Pic{\mathop{\mathrm{Pic}}\nolimits}
\def\Picardstack{\mathcal{P}\!\mathit{ic}}
\def\Picardfunctor{\mathrm{Pic}}
\def\Deformationcategory{\mathcal{D}\!\mathit{ef}}
\usepackage{cleveref}
\crefname{originaldefinition}{original definition}{original definitions}
\Crefname{originaldefinition}{Original definition}{Original definitions}
\makeatletter
\newlabel{definition-separated}{{definition-separated}{source}{}{sourcepointer-1}{}}
\newlabel{definition-separated@cref}{{[originaldefinition][0][]definition-separated}{[1][0][]source}{}{}{}}
\makeatother
\title{Four equivalent formulations of separatedness for a map of spectral spaces}
\author{The Stacks Project Authors}
\date{stacks-project:topology.tex:local-snapshot-20261012; GNU FDL 1.2 or later}
\begin{document}
\maketitle
\noindent\textit{Editorial note.} All delivered mathematical text is the author's own, copied from the pinned Stacks Project chapter source sources/stacks-topology/pinned-original-source.tex; no statement or proof step was rewritten, reordered or removed.\par\medskip
\noindent\textit{Editorial note.} Preamble provenance: the Stacks Project keeps one single shared preamble for every chapter; this card takes it directly from the staged sources/stacks-topology/preamble.tex as three verbatim spans (lines 8-12, 14-18 and 154-217), so no derived file was written and nothing was added or redefined. The wrapper class and the base packages are supplied by the assembler.\par\medskip
\noindent\textit{Editorial note.} Bibliography provenance: neither the counted statement nor its proof contains a citation (this lane rescans the whole span before writing the card), so the card carries no bibliography and no entry was reconstructed.\par\medskip
\noindent\textit{Editorial note.} Selection provenance: this card was opened under the file-based control contract (CONTROL.md), with the mechanical pre-card gate run on its statement and proof interval before the card was written.\par\medskip
\noindent\textit{Allowed standard prerequisites (editorial).}\begin{itemize}
\item The description of the diagonal of a map of spectral spaces through the topology of the fibre product. Used inside the counted proof (source lines 432-440). Turns the diagonal being a closed immersion into a condition on the diagonal as a subspace.
\end{itemize}\medskip
\ifdefined\hypertarget\hypertarget{sourcepointer-1}{}\fi
\noindent\textit{Source reference.} Original definition definition-separated: sources/stacks-topology/pinned-original-source.tex, source lines 413--413. Navigation pointer only: the counted text refers to the Stacks item definition-separated, which is an existing labeled result of this same chapter and is not part of this card. The pointer carries the target's own label as its number and its definition line as its span; its content is not reproduced here.\par\medskip
\begin{lemma}
\label{lemma-separated}
Let $f : X \to Y$ be continuous map of topological spaces.
The following are equivalent:
\begin{enumerate}
\item $f$ is separated,
\item $\Delta(X) \subset X \times_Y X$ is a closed subset,
\item given distinct points $x, x' \in X$ mapping to the same point of
$Y$, there exist disjoint open neighbourhoods of $x$ and $x'$.
\end{enumerate}
\end{lemma}

\begin{proof}
If $f$ is separated, by Definition \ref{definition-separated}, $\Delta$ is a
closed map. The fact that $X$ is closed in $X$ gives us that $\Delta(X)$ is
closed in $ X\times_Y X $. Thus (1) implies (2).

\medskip\noindent
Assume $\Delta(X) \subset X \times_Y X$ is a closed subset and let $U$
denote the complementary open. This means we have an open set
$W \subset X \times X$ such that $W \cap (X \times_Y X) = U$. However,
by definition of the product topology, if
$(x, x') \in W \cap (X \times_Y X)$,
we have $V$ and $V'$ open sets of $X$ such that $x \in V$, $x' \in V'$ and
$V \times V' \subset W$. If we had $V \cap V' \neq \emptyset$, we would have
$z \in V \cap V'$. However, $(z, z) \in X \times_Y X$, so
$(z,z) \in (V \times V') \cap (X\times_Y X) \subset U$, which is absurd.
Therefore $V \cap V' = \emptyset $, and we have two disjoint open
neighborhoods for $x$ and $x'$. It proves that (2) implies (3).

\medskip\noindent
Finally, we suppose that given distinct points $x, x' \in X$ mapping to
the same point of $Y$, there exist disjoint open neighbourhoods
of $x$ and $x'$. Let $F$ be a closed set of $X$. We will show that
$\Delta(F)$ is a closed subset of $X \times_Y X$. Let
$(x, x') \in X \times_Y X$ be a point not contained in $\Delta(F)$.
Then either $x \not = x'$ or $x \not \in F$.
In the first case, we choose disjoint open neighbourhoods
$V, V' \subset X$ of $x, x'$ and we see that $(V \times V') \cap X \times_Y X$
is an open neighbourhood of $(x, x')$ not meeting $\Delta(F)$.
In the second case, we see that $((X \setminus F) \times X) \cap X \times_Y X$
is an open neighbourhood of $(x, x')$ not meeting $\Delta(F)$.
We have shown that (3) implies (1).
\end{proof}

\end{document}
